\documentclass[11pt,a4paper]{article}

\usepackage[utf8]{inputenc}
\usepackage[cyr]{aeguill}
\usepackage[francais]{babel}

% Les packages
%\usepackage[french]{babel}
\usepackage{fancybox}
\usepackage{graphics}
\usepackage{color}
\usepackage{latexsym}
\usepackage{amsmath,amsthm,amssymb}
\usepackage[plain]{fullpage} 
\usepackage{verbatim}
\usepackage{times,epsfig}
\usepackage{listings}
\usepackage{multirow}
\usepackage{amsmath,amsthm,amssymb,amscd,txfonts,bbm} % RAJOUT ICI !!!
\usepackage{graphics,epsfig} % RAJOUT ICI !!!
%\usepackage{algorithmic}
\usepackage[final]{pdfpages} 
\usepackage{fancyvrb}
\usepackage{enumerate}
\usepackage{listingsutf8}
\usepackage{comment}
\usepackage{minted}

% small et verbatim: to be coherent with \file
\newenvironment{smallverbatim}%
{\endgraf\small\verbatim}{\endverbatim}


\newcommand{\cle}[1]{\textbf{#1}}
\def\fk#1{\textit{#1}}
\def\tble#1{\textit{#1}}

\newenvironment{qv}
{\quote\Verbatim}
{\endquote\endVerbatim}

\usepackage{lastpage}

\newcommand{\terme}[1]{\texttt{#1}}
\newcommand{\termSet}[1]{\{\texttt{#1}\}}
\newcommand{\guill}[1]{<<\,#1\,>>}
% Macro pour l'inclusion de figure  
\newcommand{\fig}[2]{%
  \begin{figure*}[hbt]
   \begin{center}
      \includegraphics[width=6cm]{#1}
  \caption{\label{#1}  #2}
  \end{center}
 \end{figure*}
}
\newcommand{\figTex}[2]{%
  \begin{figure*}[hbt]
 \centerline{\input{#1.pstex_t}}
  \caption{\label{#1}  #2}
 \end{figure*}
}

% Macro commentaires
\definecolor{turquoise}{rgb}{0.20 ,0.60, 0.60}
\definecolor{vert}{rgb}{0.40,  0.60 ,0.00}
\definecolor{violet}{rgb}{0.80 , 0.00 , 0.80}
\definecolor{bleu}{rgb}{0.20,  0.20,  0.80}
\definecolor{rose}{rgb}{1.00,  0.20,  0.40}

\lstset{basicstyle=\footnotesize,showstringspaces=false,language=XML,breaklines=false,columns=fullflexible,frame=shadowbox, captionpos=b}

\newenvironment{solution}[0]{
~\\ \color{red} \textbf{Solution : }\color{black}}{%
}
%\excludecomment{solution}


\parindent=0pt           
                             
\usepackage{url}

\newtheorem{Exercice}{Definition}


\begin{document}

\title{Sujet d'examen UE NFE204}
\author{Session FOD}
\date{8 septembre 2018}

\maketitle

\begin{center}
{\bf \Large Consignes}\\
\begin{tabular}{| p{15cm} |}
\hline
Dur\'ee de l'examen : 3h00; Documents non autoris\'es;  Calculatrice autoris\'ee\\

Les exercices sont indépendants les uns des autres. Merci de soigner votre écriture.\\
Ce sujet contient \pageref{LastPage} pages.\\

%\textit{Important}\,: nous n'évaluons pas les réponses par la syntaxe. Ne vous
%inquiétez pas d'utiliser un point-virgule à la place d'un point d'exclamation ou
%autres détails mineurs
\hline
\end{tabular}
\end{center}

%\emph{L'ensemble de l'examen doit vous permettre de montrer que vous avez assimilé les principales notions
%du cours et que vous savez les mettre en application et les expliquer. La clarté, la concision,
%la précision des réponses sera donc appréciée.}



\section*{Première partie : documents structurés (6 pts)}

Nous gérons un site de commerce électronique qui propose des services de co-voiturage.
Nous commençons avec des objectifs modestes (quelques milliers de clients) et la
base relationnelle suivante\,:

\begin{itemize}
  \item Conducteur   (\textbf{id}, nom)
  \item Client (\textbf{id}, nom)
  \item Trajet (\textbf{\textit{idConducteur}}, \textbf{\textit{idClient}}, tarif)
\end{itemize}

\vspace*{0.2cm}

Voici également un tout petit échantillon de la base.

\vspace*{0.2cm}
\begin{minipage}[b]{2.5cm}
\begin{tabular}{l|l}
\cle{id} & nom\\
\hline
1 & Albert\\
2 & Aline\\
\hline
\end{tabular}

\centerline{La table \tble{Conducteur}}
\end{minipage} \hspace*{0.5cm} \
\begin{minipage}[b]{3cm}
\begin{tabular}{l|l}
\cle{id} &  nom  \\
\hline
100 & Berthe\\
101 & Eugène \\
\hline
\end{tabular}

\centerline{La table \tble{Client}}
\end{minipage} \hspace*{0.3cm} \
\begin{minipage}[b]{3cm}
\begin{tabular}{l|c|l}
\cle{\fk{idConducteur}} & \cle{\fk{idClient}}  & tarif \\
\hline
1  & 100  & 1 \\
1 & 101  & 5 \\
2 & 101  & 3 \\
\hline
\end{tabular}

\centerline{La table \tble{Trajet}}
\end{minipage}

\vspace*{0.2cm}

On décide de remplacer la base relationnelle par une base NoSQL. 

\begin{enumerate}
  \item Donnez une première représentation de l'échantillon ci-dessus sous la forme d'un document JSON. 
      Nous l'appellerons représentation A.
  \item Proposez une autre représentation, toujours en JSON, différente de la première.  
  Nous l'appellerons représentation B.
  \item On veut calculer le tarif moyen des trajets, par conducteur (calcul C1) et par client (calcul C2).
      Indiquez pour chaque calcul la représentation (A ou B) la plus adaptée et expliquez brièvement.
     \item Est-il possible de passer de la représentation A à la représentation B par un traitement
        Map Reduce? Expliquez comment, en indiquant ce que font les fonctions de Map et de Reduce.
\end{enumerate}

\section*{Deuxième partie : recherche d'information (7 pts)}

Voici des descriptifs (abrégés) d'UEs proposées par le Cnam.

\vspace*{0.3cm}


\begin{tabular}{lp{15cm}}
     $c_1$& On étudie les systèmes d'information construits en java.\\
     $c_2$&  L'UE est consacrée aux interfaces java d'analyse de données dans les bases nosql.\\
     $c_3$& Les \textit{design patterns} en Java et leur implantation dans les \textit{frameworks} java 2.\\
     $c_4$& On montre que les systèmes nosql sont bien adaptés à l'analyse de grandes collections.\\
     $c_5$& On étudie les systèmes java en premier lieu, puis les systèmes d'analyse en second lieu\\ 
\end{tabular}

\bigskip

On va se limiter au vocabulaire (\texttt{java, système, nosql, analyse}).
(NB: on suppose
     une phase préalable de normalisation qui élimine les pluriels, majuscules, etc.)

\begin{enumerate}
   \item Donnez une matrice d'incidence contenant les tf, 
      et un tableau donnant les idf (sans le $\log$), pour les  termes précédents. Placez
      les termes en ligne, les descriptifs en colonne.

\begin{solution}
\begin{center}
\begin{tabular}{l|c|c|c|c|c}
            & \textbf{c1} & \textbf{c2} & \textbf{c3} &\textbf{c4}&\textbf{c5} \\ \hline  
   java  5/4   & 1           & 1           & 2           & 0         & 1\\
   système 5/3  & 1           & 0           & 0           & 1         & 2\\
   nosql 5/2  & 0           & 1           & 0           & 1         & 0\\
   analyse 5/3  & 0           & 1           & 0           & 1         & 1\\\hline
   \end{tabular}
\end{center}
\end{solution}

  \item Donnez les normes des vecteurs représentant ces descriptifs.
\begin{solution}
Les normes
\begin{itemize}
  \item $||c_1|| = \sqrt{1+ 1}$; $||c_2|| = \sqrt{3}$; $||c_3|| = \sqrt{4}$; 
    $||c_4|| = \sqrt{3}$, $||c_5|| = \sqrt{6}$  
\end{itemize}
\end{solution}
  
  \item Donner la similarité cosinus basée sur les tf 
  (on ignore l'idf) pour les requêtes suivantes. Expliquez brièvement la raison
  du classement obtenu.
  
  \begin{itemize}
    \item système;
    \item nosql et analyse;
    \item java et système.
  \end{itemize}
\begin{solution}
Les cosinus (requête non normalisée).

\begin{itemize}
  \item système\,: c1 : $\frac{1}{\sqrt{2}}$;  
       c4 : $\frac{1}{\sqrt{3}}$; 
       c5 : $\frac{2}{\sqrt{6}}$; 
  
  \item nosql et analyse\,: 
  
    c2 : $\frac{1+1}{\sqrt{3}}$; 
    c4 :  $\frac{1+1}{\sqrt{3}}$  
    c5 : $\frac{1}{\sqrt{6}}$; 

  \item java et système
  
    c1 : $\frac{1+1}{\sqrt{2}}$; 
    c2 : $\frac{1}{\sqrt{3}}$;
    c3 :  $\frac{2}{\sqrt{3}}$  
    c4 :  $\frac{1}{\sqrt{3}}$  
    c5 : $\frac{1+2}{\sqrt{6}}$; 
    
 \end{itemize}

\end{solution}  

   \item Pour la dernière requête, "java et système", on constate que deux
     documents sont à égalité. Est-ce que cela change 
   si on prend en compte l'idf, comment et pourquoi\,?
   
\begin{solution}
bash

système est un peu moins présent que java dans la collection, ce qui amène à classer
c4 avant c2.
\end{solution}

\item En supposant que l'on décide que deux
   termes déjà présents dans la matrice sont synonymes, comment peut-on modifier
    les informations de cette dernière (tf et idf) pour en tenir compte, sans tout recalculer\,?

\end{enumerate}


\section*{Troisème partie : systèmes distribués (4 pts)}

On vous demande d'expliquer quelques-uns des principes du \textit{consistent hashing}. 

\begin{enumerate}
   \item Expliquez comment on choisit le serveur sur lequel on place un document.
  \item Pourquoi la structure est-elle basée sur un anneau? Pourquoi pas un segment de droite par exemple?
  \item Pourquoi y a-t-il autant de valeurs sur l'anneau (2 puissance 64 typiquement)? Pourquoi ne pas commencer
      avec 8 ou 16 valeurs, et augmenter ensuite?
   \item Quelle est l'utilité de la notion de n{\oe}ud virtuel?
\end{enumerate}

\section*{Quatrième partie : questions de cours (3 pts)}

\begin{enumerate}
 \item Pourquoi faut-il stocker les paires intermédiaires sur disque dans un processus
 MapReduce? Pourquoi ne peut-on directement les transmettre aux machines qui effectuent
 la phase de Reduce?
 
\item Dans une ferme de serveur (\textit{cloud}) est-il plus efficace échanger des données
   entre les mémoires RAM de deux serveurs disctincts, ou entre la mémoire RAM et le disque d'un même serveur?
  
      \item Que pensez-vous de l'affirmation suivante\,: \emph{la réplication sert essentiellement
        à améliorer les temps de lecture}\,?
    
\end{enumerate}


\end{document}
